\documentclass{article}
\usepackage{times}
\usepackage{a4wide}
\usepackage{latexsym}
\usepackage{mathrsfs}
\usepackage{amsmath}
\usepackage{amssymb}
\usepackage{amsthm}
\usepackage{ifthen}
\usepackage{stmaryrd}
\usepackage{algorithm}
\usepackage{ctex}
%\usepackage{algorithmic}
\usepackage{algpseudocode}
\usepackage{graphics}
\usepackage{epsfig}
\usepackage{setspace}
\setstretch{1}
%\usepackage{ramacros}

\addtolength{\textwidth}{20mm}
\addtolength{\oddsidemargin}{-10mm}
\addtolength{\textheight}{18mm}
\addtolength{\topmargin}{-9mm}


\newcounter{exercise}
\setcounter{exercise}{0}
\newcommand{\exercise}{
        \addtocounter{exercise}{1}
        \vspace{0.2in}
        \noindent
        {\bf \theexercise .}
        }

\newcommand{\REMARK}[1]{}


\newcommand{\NEWPART}{\vspace{.1in}
              \noindent}

\newcommand{\<}{
    \langle}

\renewcommand{\>}{
    \rangle}

\newcommand{\ceil}[1]{\left\lceil #1 \right\rceil}

\pagestyle{plain} \pagenumbering{arabic}

\title{{\bf Assignment 2} \\ {\large Deadline: Apr. 6}}

\author{125033910296 李博洋}
\date{}

\begin{document}
\maketitle


{\large


\begin{exercise}
You are given $k$ identical eggs and an $n$ story building. You need to figure out the highest
floor $l \in {0, 1, 2, \cdots, n}$ that you can drop an egg from without breaking it. Each egg will
never break when dropped from floor $l$ or lower, and always breaks if dropped from floor
$l + 1$ or higher. ($l = $0 means the egg always breaks). Once an egg breaks, you cannot use it
any more. However, if an egg does not break, you can reuse it.

Let $f(n, k)$ be the minimum number of egg drops that are needed to find $l$ (regardless of
the value of $l$).
\end{exercise}
\\\\
\noindent\textbf{Solution:}
通过动态规划求解。还剩$j$个鸡蛋，在第$i$层楼扔鸡蛋，有以下两种情况：
\begin{itemize}
    \item 鸡蛋碎了，还剩$j-1$个鸡蛋，则需要从第$i-1$层楼开始扔鸡蛋，直到找到最高楼层。
    \item 鸡蛋没碎，还剩$j$个鸡蛋，则需要从第$i+1$层楼开始扔鸡蛋，直到找到最高楼层。
\end{itemize}
最坏情况取两者最大值，再枚举所有可以扔鸡蛋的楼层$x$取最小值，因此，有状态转移方程：
\begin{equation*}
f(i, j) = \min_{1\leq x\leq i}\Bigl(1 + \max\bigl(f(x-1, j-1), f(i-x, j)\bigr)\Bigr) \\
\end{equation*}
\noindent 注意，此处$i-x$的原因是需要检测的范围是$x+1$层楼到$i$层楼，因此需要检测的楼的数量是$i-x$层楼。

\noindent 同时，有边界条件：
\begin{gather*}
f(i, 1) = n, \\
\text{当只有一个鸡蛋时，只能逐层扔鸡蛋，直到找到最高楼层。} \\
f(n, j) = \lceil\log_2(n+1)\rceil, \quad j \geq \log_2(n+1). \\
\text{当鸡蛋数大于 $\log_2(n+1)$ 时，使用二分法最多需要 $\lceil\log_2(n+1)\rceil$ 次扔鸡蛋即可找到最高楼层。}
\end{gather*}
最终，$f(n, k)$即为所求。


\begin{exercise}
Design an efficient algorithm to work out the minimum number of characters to be deleted to make string a palindrome.
\end{exercise}
\\\\
\noindent\textbf{Solution:}
该问题可以转换为：对于字符串$s$，求出最长的回文子序列的长度，然后删除剩余的字符即可。

定义$f(i, j)$为字符串$s[i:j]$的最长回文子序列的长度，则有状态转移方程：
\begin{equation*}
f(i, j) = \begin{cases}
f(i-1, j-1) + 2 & \text{if } s[i] = s[j] \\
\max(f(i-1, j), f(i, j-1)) & \text{if } s[i] \neq s[j] \\
\end{cases}
\end{equation*}

\noindent 同时，有边界条件：
\begin{gather*}
    f(i, j) = 1, \quad i = j \\
\end{gather*}

从边界条件开始，逐步计算$f(i, j)$，最终$f(0, n-1)$即为所求（$s$的长度为$n$）。时间复杂度为$O(n^2)$。

\begin{exercise}
There are $d$ dice each having $f$ faces. Find the number of ways in which we can get sum $s$ from the faces when the dice are rolled.
\end{exercise}
\\\\
\noindent\textbf{Solution:}
定义$f(i, j)$为掷出$i$个骰子，掷出$j$的方案数，则有状态转移方程：
\begin{equation*}
    f(i, j) = \sum_{k=1}^{f} f(i-1, j-k)
\end{equation*}

\noindent 同时，有边界条件：
\begin{gather*}
    f(1, j) = 1, \quad j \leq f \\
    f(i, j) = 0, \quad j < d \text{ or } j > d*f 
\end{gather*}
从边界条件开始，逐步计算$f(i, j)$，最终$f(d, s)$即为所求。时间复杂度为$O(d*s*f)$。


\begin{exercise}
Consider the following \textbf{3-partition problem}. Given integers $a_1,\ldots,a_n$, we want to determine whether it is possible to partition of $\{1,\ldots, n\}$ into three disjoint subsets $I$, $J$, $K$ such that
$$\sum_{i\in I} a_i =\sum_{j\in J} a_j= \sum_{k\in K} a_k =\frac{1}{3} \sum_{i=1}^n a_i$$
Devise and analyze an algorithm for 3-partition that runs in time polynomial in $n$ and in $\sum_i a_i$.
\end{exercise}
\\\\
\noindent\textbf{Solution:}
    定义函数$f(i, j, k)$，$f(i, j, k)=true$ 表示从前i个元素中可以选取元素组成两个分离的子集，
使得其中一个子集元素和为$j$，另一个子集元素和为$k$，则有状态转移方程：
\begin{equation*}
    f(i, j, k) = f(i-1, j, k) \text{ or } f(i-1, j-a_i, k) \text{ or } f(i-1, j, k-a_i)
\end{equation*}

\noindent 同时，有边界条件：
\begin{gather*}
    f(i, 0, 0) = true, \quad f(0, j, k) = false, \quad f(0, 0, 0) = true
\end{gather*}

记$S = \sum_{i=1}^n a_i$，则最终答案为$f(n, S/3, S/3)$。时间复杂度为$O(n*S^2)$。



\begin{exercise}
Given an unlimited supply of coins of denominations $x_1, x_2,\dots x_n$, we wish to make change for
a value $v$ using at most $k$ coins; that is, we wish to find a set of $\leq k$ coins whose total value is $v$.
This might not be possible: for instance, if the denominations are $5$ and $10$ and $k = 6$, then we
can make change for 55 but not for 65. Give an efficient dynamicprogramming algorithm for the
following problem.
\begin{description}
\item[Input] $x_1,\ldots, x_n; k; v$.
\item[Question] Is it possible to make change for $v$ using at most $k$ coins, of denominations
$x_1,\ldots, x_n$?
\end{description}
\end{exercise}
\\\\
\noindent\textbf{Solution:}
定义函数$f(k, v)$，$f(k, v)=true$表示可以使用至多$k$个硬币组成价值为$v$的方案，则有状态转移方程：
\begin{equation*}
    f(i, j) = f(i-1, j-x_1) \text{ or } f(i-1, j-x_2) \text{ or } \cdots \text{ or } f(i-1, j-x_n)
\end{equation*}
\noindent 同时，有边界条件：
\begin{gather*}
    f(i, 0) = true, \quad f(0, v) = false, \quad v > 0, 且f(i, v) = false, \quad v < 0
\end{gather*}
最终答案为$f(k, v)$。时间复杂度为$O(n*k*v)$。


\begin{exercise}
    Given a non-empty array $n$ containing only positive integers, find if the array can be \textbf{partitioned} into two subsets such that the sum of elements in both subsets is equal.

    For example, given $n$ = [1, 5, 5, 11], The array can be partitioned as [1, 5, 5] and [11].
\end{exercise}
\\\\
\noindent\textbf{Solution:}
记$S = \sum_{i=1}^n a_i$，则问题等价于是否可以找到两个子集，使得两个子集的元素和均为$S/2$。
定义函数$f(i, j)$，$f(i, j)=true$ 表示从前$i$个元素中可以选取元素组成一个子集，使得子集的元素和为$j$，则有状态转移方程：
\begin{equation*}
    f(i, j) = f(i-1, j-a_i) \text{ or } f(i-1, j)
\end{equation*}
\noindent 同时，有边界条件：
\begin{gather*}
    f(0, j) = false, \quad f(i, 0) = ture, f(i, j) = false, \quad j < 0 \text{ 或 } i < 0
\end{gather*}
最终答案为$f(n, S/2)$。
\end{document}
